\chapter{实践卷：从 LeetCode 案例到可验证代码}

这一本书的实践卷不要求读者额外背一套模板，而是把正文中的典型题变成可运行、可对照、可扩展的代码练习。入口工程在 \filepath{books/algorithm-interview/labs/cpp20-examples}。它现在覆盖数组、双指针、滑动窗口、前缀和哈希、BFS 图搜索等基础能力，对应正文中的 LeetCode 1、3、11、53、200、209、560、994 等样本题。实践时不要只跑题目给的样例；每个题都要保留暴力版、优化版、边界用例和失败解释。

\section{实践一：数组、窗口和前缀和}

第一组练习围绕两个数组模式文件展开：

\begin{lstlisting}[language=bash]
include/algobook/array_patterns.hpp
src/array_patterns.cpp
\end{lstlisting}

它不是为了收集更多答案，而是训练从题面到不变量的过程。

\begin{enumerate}
  \item LeetCode 1 两数之和：先写双重循环暴力版，确认返回的是下标而不是数值；再写哈希表版，解释为什么必须先查 \code{target - nums[i]} 再插入当前元素，否则 \code{[3,3]} 会被错误处理。
  \item LeetCode 11 盛最多水的容器：先枚举左右边界，再推导“移动短板才可能变好”。实践代码中 \code{max_area_bruteforce} 和 \code{max_area_two_pointers} 必须对同一批输入返回相同结果。
  \item LeetCode 209 长度最小的子数组：只在数组元素全为正数时滑动窗口成立。实践时要加入含负数的反例，写清为什么窗口收缩不再安全。
  \item LeetCode 560 和为 K 的子数组：先写 \complexity{O(n^2)} 累加版，再写“当前前缀和减目标值”的哈希版。必须测试负数和全零数组，因为它们能证明窗口解法不适用。
  \item LeetCode 3 无重复字符的最长子串：暴力版保存当前集合，窗口版保存字符最后出现位置。实践重点是解释左边界为什么只能右移，不能回退。
\end{enumerate}

\begin{lstlisting}[language=bash,caption={算法实践工程的构建和测试}]
cmake -S books/algorithm-interview -B books/algorithm-interview/build/algo-debug -DCMAKE_BUILD_TYPE=Debug
cmake --build books/algorithm-interview/build/algo-debug
ctest --test-dir books/algorithm-interview/build/algo-debug --output-on-failure
\end{lstlisting}

\section{实践二：BFS 图搜索}

第二组练习在 \filepath{include/algobook/graph_patterns.hpp} 和 \filepath{src/graph_patterns.cpp}。LeetCode 994 腐烂的橘子和 LeetCode 200 岛屿数量都能用 BFS，但它们考的不是同一个细节。腐烂的橘子是多源 BFS，队列初始时要放入所有腐烂橘子；岛屿数量是连通块扫描，遇到未访问陆地就启动一次 BFS。若把这两类题都写成“套 BFS 模板”，就会漏掉初始状态、层数推进和 visited 语义。

实践报告至少回答四个问题：初始队列来自哪里，什么时候计入一层时间，访问标记什么时候写入，无法到达的节点如何被发现。对 LeetCode 994，应加入“有新鲜橘子被墙隔开”的用例；对 LeetCode 200，应加入全水、全陆地、多个岛屿相邻但不斜连的用例。

\section{实践三：用测试守住题解质量}

算法题答案进入书稿前要满足三条线。第一，至少有一个暴力或 reference 对照，能说明优化版没有改语义。第二，测试覆盖空输入、单元素、重复值、负数、极值、无答案和多答案边界。第三，复杂度解释必须来自代码结构：几层循环、每个元素入队几次、每个下标移动几次，而不是只写结论。

\begin{labbox}
扩展任务：给 LeetCode 53 最大子数组和补进实践工程。要求同时写 \complexity{O(n^2)} 固定左端累加版和 \complexity{O(n)} 动态规划版，测试全负数组、单元素、正负交替和全部为零。报告中必须从“以前一个位置结尾的最佳连续段是否值得接上当前元素”推到状态转移，不能只写公式。
\end{labbox}
